\chapter{Adder (fadd)}
\label{ch:fadd}

\section{Introduction}

\sig{fadd} computes $\mathtt{op\_a}+\mathtt{op\_b}$. Two \sig{unpack} modules
split the operands. The \sig{add} module aligns the mantissas and adds them in
three pipeline stages. \sig{norm} normalizes the sum, and \sig{pack} rounds it
and forms the result, as Figure~\ref{fig:fadd} shows.

\begin{figure}[H]
\centering
\begin{tikzpicture}[node distance=7mm]
  \node[ugblock] (u) {\sig{unpack}\\(2)};
  \node[ugblock,right=of u] (a) {\sig{add}\\3 cycles};
  \node[ugblock,right=of a] (n) {\sig{norm}\\1 cycle};
  \node[ugblock,right=of n] (p) {\sig{pack}\\1 cycle};
  \draw[ugarrow] ($(u.west)-(9mm,0)$) node[above right,font=\footnotesize\ttfamily,xshift=-1mm] {op} -- (u);
  \draw[ugarrow] (u) -- (a); \draw[ugarrow] (a) -- (n); \draw[ugarrow] (n) -- (p);
  \draw[ugarrow] (p) -- ($(p.east)+(9mm,0)$) node[above left,font=\footnotesize\ttfamily,xshift=1mm] {res};
\end{tikzpicture}
\caption{\sig{fadd} Block Diagram}
\label{fig:fadd}
\end{figure}

\section{Features}

\begin{itemize}
\item Addition with a latency of 5 clock cycles.
\item One new operation per clock cycle.
\item Subnormal operands and results.
\item Round to nearest, ties to even.
\item NaN and infinity handling, selected by \sig{SPECIAL\_CASES}.
\end{itemize}

\section{Parameters}

\sig{fadd} has the parameters of Table~\ref{tab:common_params}.

\section{Interface Signals}

\sig{fadd} has the signals of Table~\ref{tab:common_sigs}.

\Needspace*{23\baselineskip}
\section{Instantiation}

The example below instantiates \sig{fadd} for binary32. Compile it with the files
of Table~\ref{tab:files}.

\begin{vcode}
fadd #(
   .DATA_W       (32),
   .EXP_W        (8),
   .SPECIAL_CASES(1)
) fadd0 (
   .clk    (clk),
   .rst    (rst),
   .valid_i(valid_i),
   .valid_o(valid_o),
   .op_a   (op_a),
   .op_b   (op_b),
   .res    (res)
);
\end{vcode}

\section{Operation}

\sig{fadd} follows the protocol of Section~\ref{sec:protocol}, with a latency of
5 clock cycles. The unit samples the operands in the \sig{valid\_i} cycle, so they
can change in the next cycle. A new operation can start in every clock cycle, and
the results come out in order.

Table~\ref{tab:fadd_special} lists the results for special operands when
\sig{SPECIAL\_CASES} is 1. The sum of two zeros of the same sign keeps that sign.
An exact cancellation gives $+0$.

\begin{twotable}{\sig{fadd} Special Results}{tab:fadd_special}{Operands}{Result}{7cm}
Either operand is NaN & NaN \\ \hline
$+\infty$ and $-\infty$ & NaN \\ \hline
Infinity and a finite number, or two infinities of the same sign & That infinity \\
\end{twotable}
